\chapter{Положительность скрученного следа и модулярного состояния при обмене}

Вещественно-скалярный обмен задаёт допустимую скрученную геометрию, но
обычный спектральный след сохраняет неверный смешанный знак. Последняя
возможность этого направления --- положительный точный \(\rho\)-след,
функционал KMS или модулярный вес. Их существование проверяется до записи
взвешенного спектрального действия.

\section{Абстрактное препятствие для скрученного следа}

Пусть \(e_+,e_-\) --- примитивные центральные идемпотенты с
\begin{equation}
 \rho(e_+)=e_-,\qquad\rho(e_-)=e_+,\qquad e_+e_-=0.
\end{equation}
Правый \(\rho\)-след удовлетворяет
\begin{equation}
 \tau_\rho(ab)=\tau_\rho(\rho(b)a).
 \label{eq:v5-rho-trace-definition}
\end{equation}
При \(a=b=e_+\) получаем
\(\tau_\rho(e_+)=\tau_\rho(e_-e_+)=0\); аналогично
\(\tau_\rho(e_-)=0\). Следовательно, вся удвоенная вещественная компонента
аннулируется.

\begin{theorem}[Невозможность точного следа при обмене центральных идемпотентов]
На унитальной алгебре с двумя ортогональными центральными идемпотентами,
которые переставляются автоморфизмом \(\rho\), не существует точного
линейного \(\rho\)-следа. Положительность для этого вывода не требуется.
\end{theorem}

Положительные скрученные следы существуют в иных ситуациях, когда
скручивание реализуется положительным модулярным оператором; см.
\path{arXiv:1607.03834} и \path{arXiv:1111.6328}. Внешний обмен центральных
идемпотентов устроен иначе.

\section{Препятствие для операторного веса}

Пусть \(\operatorname{Tr}_W(T)=\Tr(WT)\) реализует скручивание. Невырожденность
матричного следа требует
\begin{equation}
 YW=W\rho(Y).
 \label{eq:v5-weight-intertwiner}
\end{equation}
Для \(P_+=\pi(e_+)\), \(P_-=\pi(e_-)\) в явном 18-мерном представлении
\begin{equation}\rank P_+=6,\qquad\rank P_-=3.\end{equation}
Вес должен обменивать подпространства размерностей шесть и три. На их сумме
его ранг не превосходит \(2\min(6,3)=6\), а дополнительное неподвижное
подпространство имеет размерность девять. Поэтому
\begin{equation}\rank W\le15<18.\end{equation}
Граница достижима, но всякий такой оператор имеет ядро. Положительного
обратимого веса не существует.

\section{Препятствие для конечномерного модулярного состояния}

Для точного конечномерного состояния с матрицей плотности \(K>0\)
модулярная эволюция является внутренней:
\begin{equation}\sigma_t(X)=K^{it}XK^{-it}.\end{equation}
Внутренний автоморфизм поточечно сохраняет центр и не может обменять
проекторы разных рангов. Следовательно, вещественно-скалярный обмен не
является модулярным автоморфизмом точного конечномерного состояния. См.
конечномерную реализацию Томиты--Такесаки в \path{arXiv:1301.1836}.

Переход к алгебре типа III не исключён, но он заменяет конечную спектральную
архитектуру, а не завершает её.

\section{Следствие для нормировки}

Абстрактный \(\rho\)-след не нормирует кинетический сектор обмена, а
операторный вес всегда вырожден. Нормировка на оставшейся опоре потребовала
бы нового проектора или псевдообратного оператора, то есть дополнительных
данных меры. Неопределённый знакопеременный функционал также не является
положительной евклидовой мерой.

\begin{nogocriterion}[Вырожденный скрученный вес]
Неточный \(\rho\)-след, вырожденный вес, неопределённый оператор или замена
на алгебру типа III не считаются замыканием конечного скручивания: каждый
вариант меняет класс меры или добавляет невыведенные данные опоры.
\end{nogocriterion}

\section{Вывод}

\begin{protocol}
\begin{enumerate}
 \item точный абстрактный \(\rho\)-след: \textbf{невозможен};
 \item положительный обратимый 18-мерный вес: \textbf{невозможен};
 \item точная конечномерная модулярная реализация обмена:
 \textbf{невозможна};
 \item общая нормировка кинетики и кривизны: \textbf{не получена};
 \item конечное вещественно-скалярное скручивание:
 \textbf{направление закрыто};
 \item тип III или неопределённая мера: \textbf{новая архитектура};
 \item физическое замыкание: \textbf{не достигнуто}.
\end{enumerate}
\end{protocol}

Следующий оставшийся класс --- производная геометрия и BV--BFV. Проверка
\path{version5_derived_moment_map_minimal_data_gate} должна установить
минимальные градуированные данные, превращающие отображение момента в
дифференциал или главное уравнение, не вводя нужный квадрат как новую
аксиому.

Машинный сертификат сохранён в
\path{s2t_v5_flip_twisted_trace_positivity_gate_results.json}.